Drug--target affinity (DTA) prediction estimates the strength of interaction between a small molecule and a protein target. This project implements a Davis-focused MGraphDTA model that accepts a drug SMILES string and a protein amino-acid sequence. RDKit converts each molecule into a graph with 22-dimensional atom features and bidirected bond connectivity. A multiscale graph neural network (MGNN) produces a 96-dimensional molecular embedding. The target sequence is tokenized, padded or truncated to 1,200 positions, embedded at 128 dimensions, and encoded by three convolution branches in a multiscale convolutional neural network (MCNN) to produce a 96-dimensional protein embedding. Their concatenation is passed to a multilayer perceptron (MLP) that predicts a continuous Davis affinity value.

The primary verified run used one random seed, a 90/10 training/validation split within the provided training fold, and an independent test fold that is loaded after checkpoint selection. Its best validation epoch was 91 (MSE 0.2323628495). Independent test metrics were MSE 0.2570411053, RMSE 0.5069922142, concordance index (CI) 0.8790452402, and modified squared correlation coefficient ($r_m^2$) 0.6817055903. These results belong to this project protocol and are not an exact reproduction of the original paper's evaluation. The project also implements atom-level Grad-AAM visualizations, an experimental sequence-adjacency protein GNN (test MSE 0.651731), and five separate experimental protein--ligand structure demonstrations. These additions support analysis and explanation; they do not establish experimentally validated binding sites. A preliminary RCSB structure audit remains incomplete. The current automated suite contains 36 passing tests; an earlier project milestone reported 31. The RCSB aggregate counts are project-reported but cannot be independently recomputed because the audit CSV/JSON is unavailable.

